\documentclass[11pt]{article}
\input{../_estilo.tex}
\renewcommand{\tituloejercicio}{Tema 1 --- Ejercicio 3}

\begin{document}

{\LARGE\bfseries\color{azultitulo} Ejercicio 3 --- Una pila con \texttt{minimo} en $O(1)$}

\vspace{2pt}
\textbf{Dificultad:} $\bigstar$ \quad \textbf{Conferencias:} C3 (Correctitud)
\vspace{8pt}
\hrule
\vspace{12pt}

\section*{Enunciado}
\begin{enunciado}
Una pila con una operación \texttt{minimo} que cuesta $O(1)$ en el peor caso. Diseña la estructura,
enuncia el invariante que la hace correcta y da el costo de cada operación.
\end{enunciado}

\section{El problema}

Una pila soporta \texttt{push(x)}, \texttt{pop()} y ahora también \texttt{minimo()}, que debe devolver el
menor elemento actualmente en la pila. La solución ingenua --- recorrer la pila entera cada vez que se pide
el mínimo --- cuesta $O(n)$ por consulta. Queremos las tres operaciones en $O(1)$ peor caso, sin
amortizar nada (a diferencia de la conferencia 4, aquí \emph{cada} operación individual tiene que ser
barata, no solo el promedio de una secuencia).

\textbf{Especificación (C3 §1).} Precondición: \texttt{pop()} y \texttt{minimo()} se llaman solo sobre una
pila no vacía (igual que en cualquier pila estándar, donde \texttt{pop} de una pila vacía no está
definido). Postcondición: \texttt{minimo()} devuelve el menor elemento entre los que están actualmente
apilados, sin modificar la pila.

\section{Diseño: dos pilas paralelas}

La idea es mantener, junto a la pila principal, una segunda pila que en todo momento tiene en su tope el
mínimo de lo que hay apilado hasta ahí.

El código real y ejecutable vive en \texttt{verif.py}, en esta carpeta; aquí va el pseudocódigo:

\begin{lstlisting}[style=pseudo]
estructura PilaMin:
    principal (*$\leftarrow$*) pila vacia
    minimos   (*$\leftarrow$*) pila vacia      // minimos.cima = minimo de todo lo que hay en principal

    funcion push(x):
        principal.apilar(x)
        si minimos vacia o x (*$\le$*) minimos.cima:
            minimos.apilar(x)
        sino:
            minimos.apilar(minimos.cima)        // repite el minimo anterior

    funcion pop():
        minimos.desapilar()
        devolver principal.desapilar()

    funcion minimo():
        devolver minimos.cima
\end{lstlisting}

La tabla siguiente traza un ejemplo concreto, para que se vea cómo las dos pilas evolucionan juntas (no
es la demostración --- esa es la sección 3 --- solo una corrida ilustrativa):

\begin{center}
\begin{tabular}{@{}lll@{}}
\toprule
\textbf{Operación} & \textbf{principal} & \textbf{minimos} \\
\midrule
\texttt{push(5)} & $[5]$ & $[5]$ \\
\texttt{push(3)} & $[5,3]$ & $[5,3]$ \\
\texttt{push(4)} & $[5,3,4]$ & $[5,3,3]$ \\
\texttt{push(1)} & $[5,3,4,1]$ & $[5,3,3,1]$ \\
\texttt{pop()} & $[5,3,4]$ & $[5,3,3]$ \\
\texttt{push(2)} & $[5,3,4,2]$ & $[5,3,3,2]$ \\
\bottomrule
\end{tabular}
\end{center}

En cada fila, el último valor de \texttt{minimos} es exactamente el mínimo de todo \texttt{principal} en
ese momento --- es el invariante de la sección siguiente, en acción.

\section{Invariante}

\begin{resultado}{Invariante}
Después de cualquier secuencia de operaciones, \texttt{principal} y \texttt{minimos} tienen la misma
longitud y, para todo $i$: \texttt{minimos[i]} $=\min(\texttt{principal[0..i]})$.
\end{resultado}

A diferencia del invariante de un ciclo (que se demuestra sobre el índice de la vuelta), este es un
invariante \textbf{sobre la secuencia de operaciones}: se demuestra por inducción en el número de
operaciones realizadas. Es del mismo tipo que los dos invariantes de Gale-Shapley de C3 §7 (``un receptor
nunca vuelve a quedar libre'', ``un proponente solo empeora''), que también son afirmaciones sobre el
estado después de cualquier prefijo de ejecución, no sobre un ciclo con contador.

\textbf{Inicialización.} Pila vacía: ambas listas vacías (misma longitud, 0), la afirmación es vacuamente
cierta.

\textbf{Mantenimiento.} Supongamos el invariante cierto antes de una operación.
\begin{itemize}
\item \texttt{push(x)}: las posiciones $0,\dots,i{-}1$ de ambas pilas no se tocan, así que siguen cumpliendo
el invariante por hipótesis, y ambas pilas crecen una posición, así que siguen con la misma longitud.
Falta verificar la posición nueva, $i$ (donde queda $x$). Si $i=0$ la pila estaba vacía: el código guarda $x$ en
\texttt{minimos[0]}, y $\min(\texttt{principal[0..0]})=x$. Si $i\ge 1$, por hipótesis,
\texttt{minimos[i-1]} $=\min(\texttt{principal[0..i-1]})$. El código compara $x$ con ese valor y guarda el
menor de los dos en \texttt{minimos[i]}, que es exactamente
$\min(\texttt{principal[0..i-1]}\cup\{x\}) = \min(\texttt{principal[0..i]})$.
\item \texttt{pop()}: se elimina la última posición de ambas pilas a la vez (siguen con la misma longitud). El
invariante para las posiciones restantes $0,\dots,i{-}1$ no cambió (no se tocaron esos valores), así que sigue valiendo.
\end{itemize}

\textbf{Conclusión.} En cualquier momento con la pila no vacía (la precondición de la sección 1),
por el invariante aplicado al último índice, \texttt{minimos[-1]} $=\min(\texttt{principal})$, que es
justo lo que \texttt{minimo()} devuelve. Esto da la \textbf{correctitud parcial}.

\textbf{Terminación.} A diferencia de un algoritmo con un ciclo (donde hace falta un potencial que baje),
aquí no hay nada que demostrar: \texttt{push}, \texttt{pop} y \texttt{minimo} son código lineal, sin
ciclos ni recursión, así que cada uno termina trivialmente en una cantidad fija de pasos. Correctitud
parcial más terminación dan la \textbf{correctitud total} (C3 §1).

\section{Costo}

\begin{center}
\begin{tabular}{@{}lll@{}}
\toprule
\textbf{Operación} & \textbf{Tiempo} & \textbf{Por qué} \\
\midrule
\texttt{push} & $O(1)$ & una comparación y dos operaciones de apilado \\
\texttt{pop} & $O(1)$ & dos operaciones de desapilado del tope \\
\texttt{minimo} & $O(1)$ & un acceso al tope \\
\bottomrule
\end{tabular}
\end{center}

\textbf{Modelo para el peor caso.} La garantía estricta de $O(1)$ supone que las pilas subyacentes implementan
apilar, desapilar y consultar el tope en $O(1)$ peor caso; por ejemplo, mediante nodos enlazados, o con un
arreglo cuya capacidad ya esté garantizada y que no se redimensione durante la operación. \texttt{verif.py}
usa listas de Python solo para comprobar la correctitud funcional: en un arreglo dinámico, \texttt{append} y una
posible reducción de capacidad pueden ser amortizados, así que esa implementación de prueba no demuestra por sí
sola la cota de peor caso.

\textbf{Memoria:} $\Theta(n)$ adicional en el peor caso, con $n$ el número de elementos actualmente
apilados --- la segunda pila \texttt{minimos}, del mismo tamaño que la pila principal.

\textbf{Verificación.} La demostración de la sección 3 ya garantiza que \texttt{minimo()} es correcto; se
corrió de todas formas como chequeo empírico adicional (en \texttt{verif.py}), comparando contra un oráculo
de fuerza bruta (\texttt{min(lista)} sobre una lista de Python que refleja el estado real) sobre 2000
secuencias aleatorias de \texttt{push}/\texttt{pop}, consultando \texttt{minimo()} después de cada operación:

\begin{lstlisting}[style=salida]
PilaMin (dos pilas): 2000 secuencias aleatorias verificadas, sin fallos
\end{lstlisting}

\begin{resultado}{Respuesta a la pregunta del enunciado}
\begin{itemize}[leftmargin=14pt]
\item \textbf{Estructura:} dos pilas paralelas, \texttt{principal} y \texttt{minimos}, donde
\texttt{minimos} guarda en cada posición el mínimo de \texttt{principal} hasta ahí (sección 2).
\item \textbf{Invariante:} ambas pilas de la misma longitud y \texttt{minimos[i]}
$=\min(\texttt{principal[0..i]})$ para todo $i$, demostrado por inducción sobre la secuencia de operaciones, no sobre un índice de ciclo (sección 3).
\item \textbf{Costo:} \texttt{push}, \texttt{pop} y \texttt{minimo} son los tres $O(1)$ peor caso (no
amortizado, bajo el modelo de implementación de la sección 4); memoria $\Theta(n)$ adicional (sección 4).
\end{itemize}
\end{resultado}

\section*{Errores que hay que vigilar}

\begin{alerta}{}
\begin{itemize}[leftmargin=14pt]
\item \textbf{Usar \texttt{<} en vez de \texttt{<=} al comparar \texttt{x} con el mínimo actual en
\texttt{push}.} Con \texttt{<} estricto, cuando se apila un valor \textbf{igual} al mínimo actual, el
código no lo agregaría de nuevo a \texttt{minimos} con su propia entrada --- pero como aquí \emph{siempre}
se agrega una entrada por cada \texttt{push} (se repite el valor anterior si no hay nuevo mínimo), este
error concreto no rompe esta implementación particular. El error real aparece en variantes que
\textbf{solo} apilan en \texttt{minimos} cuando hay un mínimo \emph{estricto} (para ahorrar memoria,
guardando pares \texttt{(valor, repeticiones)} en vez de una entrada por cada \texttt{push}): ahí, si se
usa \texttt{<} en vez de \texttt{<=} al decidir si el valor nuevo iguala al mínimo vigente, un
\texttt{pop()} puede descontar una repetición de más y sacar el mínimo de la pila de mínimos cuando en
realidad todavía queda una copia de ese valor en la pila principal. El arreglo es tratar la igualdad como
``suma una repetición más'', no como ``ignóralo''.
\item \textbf{Olvidar que \texttt{pop()} tiene que sacar de \emph{las dos} pilas a la vez.} Si solo se saca
de \texttt{principal}, el invariante se rompe de inmediato: \texttt{minimos} queda con una entrada de más,
desalineada con \texttt{principal}, y \texttt{minimo()} empieza a devolver el mínimo de un conjunto que ya
no es el que está apilado.
\item \textbf{Pensar que una sola pila de este tipo resuelve ``mínimo de una cola''.} Con una pila, deshacer la última operación
(\texttt{pop}) es barato porque es simétrico a \texttt{push}. Con una \textbf{cola} (FIFO) no hay esa
simetría, y el problema de mantener el mínimo (o el máximo) con \texttt{encolar}/\texttt{desencolar} en
$O(1)$ amortizado es genuinamente más difícil --- se resuelve con una \textbf{cola doble monótona} (un
\texttt{deque} donde se van descartando desde atrás los elementos que nunca podrán ser el mínimo porque hay
uno más chico y más reciente), o con dos pilas \texttt{PilaMin} haciendo de cola (la de C4 §2), también
$O(1)$ amortizado.
\end{itemize}
\end{alerta}

\end{document}
